\begin{frame}
  \frametitle{Weak Form}
  Note the following vector calculus identity for a scalar function $\psi$ and vector $\vec{A}$:
  \begin{equation}
    \label{eq:vec-calc-identity}
    \div \left(\psi\vec{A}\right)=\psi\div\vec{A}+\vec{A}\cdot\grad\psi
  \end{equation}
  Let $v=v(x,y,\theta,\varphi)$ be a test function.
  Multiplying \cref{eq:even-parity-equation} by $v$, integrating over an angular and spatial element gives, and applying \cref{eq:vec-calc-identity} twice gives the weak form:
  \begin{multline}\label{eq:general-weak-form}
    -\frac{1}{\sigt}\int_{\mathcal{A}}\oint_{\partial \tau}v\left(\ohat\cdot\grad{\psiplus}\right)\ohat\cdot\hat{n}\dd s\dd\ohat+\frac{1}{\sigt}\intelems\left(\ohat\cdot\grad{\psiplus}\right)\left(\ohat\cdot\grad{v}\right)\elemdifs+\\+
    \intelems v\sigt \psiplus\elemdifs=\intelems vS\elemdifs
  \end{multline}
  where $\int_\tau$ is the integral over a triangular spatial element, $\int_\mathcal{A}$ is the integral over an angular element, and $\hat{n}$ is the unit normal vector on in the spatial domain.
\end{frame}

\begin{frame}
  \frametitle{Boundary Conditions}
  Per Lewis \& Miller~\cite{lewisComputationalMethodsNeutron1993}, the vacuum boundary conditions are:
  \begin{equation}\label{eq:vacuum-bc}
    \ohat\cdot\grad{\psiplus}\pm\sigt\psiplus=0,\quad x,y\in \partial\Gamma_v,\ohat\cdot\hat{n}\lessgtr 0
  \end{equation}
  while the reflective boundary conditions are:
  \begin{equation}\label{eq:reflect-bc}
    \psiplus\left(\ohat\right)=\psiplus\left(\ohat '\right)
  \end{equation}
  \Cref{eq:reflect-bc} is trivially satisfied by the weak form, so it will not need further discussion.
  Inserting \cref{eq:vacuum-bc} into \cref{eq:general-weak-form}'s boundary condition gives:
  \begin{equation}
    -\frac{1}{\sigt}\int_{\mathcal{A}}\oint_{\partial \tau}v\left(\ohat\cdot\grad{\psiplus}\right)\ohat\cdot\hat{n}\dd s\dd\ohat=\pm \int_\mathcal{A}\oint_{\partial\tau_v}v\psiplus\left(\ohat\cdot\hat{n}\right)\dd s\dd \ohat
  \end{equation}
  Notice that the sign in the boundary term is guaranteed to be positive.

\end{frame}

\begin{frame}
  \frametitle{Weak Form}
  Noting the following definitions for $\theta\in[0,\pi]$, $\varphi\in[0,\pi/2]$:
  \begin{equation}
    \ohat\cdot\grad{\psiplus}=\mu\pdv{\psiplus}{x}+\eta\pdv{\psiplus}{y}+\xi\pdv{\psiplus}{z}
  \end{equation}
  \begin{equation}
    \ohat=\left(\underbrace{\cos\varphi}_\mu,\underbrace{\cos\theta\sin\varphi}_\eta,\underbrace{\sin\theta\sin\varphi}_\xi\right)
  \end{equation}
  and disregarding the $z-$ component (for 2D), we finally arrive at:
  \begin{multline}\label{eq:weak-form}
    \pm\int_{\mathcal{A}}\oint_{\partial\tau}v\psiplus\left(\mu n_x+\eta n_y\right)\dd s\dd\ohat +\\+\frac{1}{\sigt}
    \intelems\left(\mu\pdv{\psiplus}{x}+\eta\pdv{\psiplus}{y}\right)\left(\mu\pdv{v}{x}+\eta\pdv{v}{y}\right)\elemdifs+\\
    +\sigt\intelems v \psiplus\elemdifs=\intelems vS\elemdifs
  \end{multline}
\end{frame}

\begin{frame}
  \frametitle{Function Approximation}
  Let
  \begin{equation}
    \label{eq:function-approx}
    \psiplus \left(x,y,\theta,\varphi\right)\approx\sum_i\sum_j \psi_{ij}T_i(x,y)A_j(\theta,\varphi)
  \end{equation}
  for $i\in[1,n_t]$ for $n_t$ triangles in space, $j\in[1,n_a]$ for $n_a$ elements in solid angle.
  Inserting \cref{eq:function-approx} into \cref{eq:weak-form} and combining summations gives:
  \begin{multline}\label{eq:function-exp-fem}
    \pm\int_{\mathcal{A}}\oint_{\partial\tau}v\sum_{i,j} \psi_{ij}T_iA_j\left(\mu n_x+\eta n_y\right)\dd s\dd\ohat +\\+
    \frac{1}{\sigt}\intelems \left(\sum_{i,j}\left(\mu\pdv{T_i}{x}+\eta\pdv{T_i}{y}\right)A_j\psi_{ij}\right)\left(\mu\pdv{v}{x}+\eta\pdv{v}{y}\right)\elemdifs+\\+
    \sigt\intelems v\sum_{i,j} T_i A_j \psi_{ij}\elemdifs = \intelems vS\elemdifs
  \end{multline}
\end{frame}

\begin{frame}
  \frametitle{Galerkin Method}
  For the test function, we choose the same basis functions as in \cref{eq:function-approx}:
  \begin{equation}\label{eq:galerkin-test-function}
    v(x,y,\theta,\varphi)=\sum_m\sum_n T_m(x,y)A_n(\theta,\varphi)
  \end{equation}
  where $m\in[1,n_t]$, $n\in[1,n_a]$.
  Inserting \cref{eq:galerkin-test-function} into \cref{eq:function-exp-fem} gives:
  \begin{multline}
    \pm\int_\mathcal{A}\oint_{\partial \tau}\sum_i\sum_j\sum_m\sum_n A_jA_n T_i T_m\left(\mu n_x+\eta n_y\right)\dd s\dd\ohat\psi_{ij}+\\+
    \frac{1}{\sigt}\intelems\sum_i\sum_j\sum_m\sum_n\left(\mu^2\pdv{T_i}{x}\pdv{T_m}{x}+\eta^2\pdv{T_i}{y}\pdv{T_m}{y}+\mu\eta\pdv{T_i}{y}\pdv{T_m}{x}+\mu\eta\pdv{T_i}{x}\pdv{T_m}{y}\right)\times\\\times
    A_jA_n\psi_{ij}+\sigt\int_{\mathcal{A}}\sum_n\sum_jA_nA_j\dd\ohat\int_\tau\sum_m\sum_iT_mT_i\dd\tau\psi_{ij}=\\=
    \int_\mathcal{A}\sum_nA_n\dd \ohat\int_\tau\sum_mT_mS\dd\tau
  \end{multline}
\end{frame}

\begin{frame}
  \frametitle{Galerkin Method}
  Or, more compactly:
  \begin{multline}\label{eq:compact-fem}
    \pm\int_\mathcal{A}\oint_{\partial \tau}\sum_{i,j,m,n}A_jA_n T_i T_m\left(\mu n_x+\eta n_y\right)\dd s\dd\ohat\psi_{ij}+\\+
    \frac{1}{\sigt}\intelems\sum_{i,j,m,n}\left(\mu^2\pdv{T_i}{x}\pdv{T_m}{x}+\eta^2\pdv{T_i}{y}\pdv{T_m}{y}+\mu\eta\pdv{T_i}{y}\pdv{T_m}{x}+\mu\eta\pdv{T_i}{x}\pdv{T_m}{y}\right)\times \\\times A_jA_n\psi_{ij}+
    \sigt\int_\mathcal{A}\sum_{j,n}A_jA_n\dd\ohat\int_\tau\sum_{i,m}T_mT_i\dd\tau\psi_{ij}=\\=
    \int_\mathcal{A}\sum_nA_n\dd \ohat\int_\tau\sum_mT_mS\dd\tau
  \end{multline}
\end{frame}

\begin{frame}
  \frametitle{Source Term}
  The source term, $S$, up until this point has been left rather ambiguous.
  Inserting the definition in \cref{eq:fully-simplified-bte} into the source integral term from \cref{eq:compact-fem}:
  \begin{multline}
    \int_\mathcal{A}\sum_nA_n\dd \ohat\int_\tau\sum_mT_mS\dd\tau =
    \frac{1}{4\pi}\sum_{i,j,m,n}\left[\Sigma_{s,i}\int_\tau T_iT_m\dd\tau\int_\mathcal{A} A_j\dd\ohat \int_\mathcal{A} A_n\dd\ohat \psi_{ij}\right]+\\+
    \frac{1}{4\pi}\sum_{m,n} \left[Q_m\int_\tau  T_m\dd\tau\int_\mathcal{A}A_n\dd\ohat\right]
  \end{multline}
  where we have approximated the scattering cross section and volumetric source to be constant within the spatial element.
\end{frame}

\begin{frame}[allowframebreaks]
  \frametitle{Kronecker Product}
  Want to rewrite \cref{eq:compact-fem} to leverage Kronecker products.
  The streaming term requires some rearranging.
  Beginning with the boundary term:
  \begin{subequations}
    \begin{multline}
      \pm\int_\mathcal{A}\oint_{\partial \tau}\sum_{i,j,m,n}A_jA_n T_i T_m\left(\mu n_x+\eta n_y\right)\dd s\dd\ohat \psi_{ij}= \\ =
      \pm\sum_{i,j,m,n}\left[\int_\mathcal{A}\mu A_jA_n\dd\ohat\oint_{\partial \tau_v}T_iT_m n_x\dd s\right]\psi_{ij}+\\
      \pm\sum_{i,j,m,n}\left[\int_\mathcal{A}\eta A_jA_n\dd\ohat\oint_{\partial \tau_v}T_iT_m n_y\dd s\right]\psi_{ij}
    \end{multline}
    \framebreak

    Now the cell term:
    \begin{multline}
      \frac{1}{\sigt}\intelems\sum_{i,j,m,n}\left(\mu^2\pdv{T_i}{x}\pdv{T_m}{x}+\eta^2\pdv{T_i}{y}\pdv{T_m}{y}+\mu\eta\pdv{T_i}{y}\pdv{T_m}{x}+\mu\eta\pdv{T_i}{x}\pdv{T_m}{y}\right)\times \\\times A_jA_n\psi_{ij} =\\=
      \frac{1}{\sigt}\sum_{i,j,m,n}\left[\int_\mathcal{A}\mu^2A_jA_n\dd \ohat\int_\tau\pdv{T_i}{x}\pdv{T_m}{x}\dd\tau\right]\psi_{ij}+\\+
      \frac{1}{\sigt}\sum_{i,j,m,n}\left[\int_\mathcal{A}\eta^2A_jA_n\dd \ohat\int_\tau\pdv{T_i}{y}\pdv{T_m}{y}\dd\tau\right]\psi_{ij}+\\+
      \frac{1}{\sigt}\sum_{i,j,m,n}\left[\int_\mathcal{A}\mu\eta A_jA_n\dd \ohat\int_\tau\pdv{T_i}{x}\pdv{T_m}{y}\dd\tau\right]\psi_{ij}+\\+
      \frac{1}{\sigt}\sum_{i,j,m,n}\left[\int_\mathcal{A}\mu\eta A_jA_n\dd \ohat\int_\tau\pdv{T_i}{y}\pdv{T_m}{x}\dd\tau\right]\psi_{ij}
    \end{multline}
  \end{subequations}

  \framebreak
  Let:
  \begin{subequations}
    \begin{equation}
      \left[K_{\mu\mu}\right]_{jn}=\int_\mathcal{A}\mu^2A_jA_n\dd \ohat
    \end{equation}
    \begin{equation}
      \left[K_{\eta\eta}\right]_{jn}=\int_\mathcal{A}\eta^2A_jA_n\dd \ohat
    \end{equation}
    \begin{equation}
      \left[K_{\mu\eta}\right]_{jn}=\int_\mathcal{A}\mu\eta A_jA_n\dd \ohat
    \end{equation}
  \end{subequations}
  \framebreak

  then:
  \begin{subequations}
    \begin{equation}
      \left[M_{xx}\right]_{im} =\int_\tau\pdv{T_i}{x}\pdv{T_m}{x}\dd\tau
    \end{equation}
    \begin{equation}
      \left[M_{yy}\right]_{im} =\int_\tau\pdv{T_i}{y}\pdv{T_m}{y}\dd\tau
    \end{equation}
    \begin{equation}
      \left[M_{xy}\right]_{im} =\int_\tau\pdv{T_i}{x}\pdv{T_m}{y}\dd\tau
    \end{equation}
    \begin{equation}
      \left[M_{yx}\right]_{im} =\int_\tau\pdv{T_i}{y}\pdv{T_m}{x}\dd\tau
    \end{equation}
  \end{subequations}
  \framebreak

  with:
  \begin{subequations}
    \begin{equation}
      \left[B_{x}\right]_{im}=\oint_{\partial \tau_v}T_iT_m n_x\dd s
    \end{equation}
    \begin{equation}
      \left[B_{y}\right]_{im}=\oint_{\partial \tau_v}T_iT_m n_y\dd s
    \end{equation}
    \begin{equation}
      \left[B_\mu\right]_{jn}=\int_\mathcal{A}\mu A_jA_n\dd\ohat
    \end{equation}
    \begin{equation}
      \left[B_\eta\right]_{jn}=\int_\mathcal{A}\eta A_jA_n\dd\ohat
    \end{equation}
  \end{subequations}
  \framebreak

  and finally:
  \begin{subequations}
    \begin{equation}
      \left[K\right]_{jn}=\int_\mathcal{A}A_jA_n\dd\ohat
    \end{equation}
    \begin{equation}
      \left[M\right]_{im}=\int_\tau T_mT_i\dd\tau
    \end{equation}
    \begin{equation}
      \left[S\right]_{im}=\frac{\Sigma_{s,i}}{4\pi}\int_\tau T_iT_m\dd\tau
    \end{equation}
    \begin{equation}
      \vec{Q}_m=\frac{Q_m}{4\pi}\int_\tau T_m\dd\tau
    \end{equation}
    \begin{equation}
      \vec{A}_n=\int_\mathcal{A}A_n\dd\ohat
    \end{equation}
    \begin{equation}
      \vec{C}_j=\int_\mathcal{A}A_j\dd\ohat
    \end{equation}
    \begin{equation}
      \left[\Psi\right]_{ij}=\psi_{ij}
    \end{equation}
  \end{subequations}
  \framebreak

  Combining these all into \cref{eq:compact-fem} gives:
  \begin{multline}
    \left[B_\mu\otimes B_x+B_\eta\otimes B_y\right]\Psi +\\+
    \frac{1}{\sigt}\left[K_{\mu\mu}\otimes M_{xx}+K_{\eta\eta}\otimes M_{yy}+K_{\mu\eta}\otimes M_{xy}+K_{\mu\eta}\otimes M_{yx}\right]\Psi+\\+
    \sigt K\otimes M\Psi = \left[S\otimes \left(\vec{C}\otimes\vec{A}\right)\right]\Psi+\vec{Q}\otimes \vec{A}
  \end{multline}
  where $\otimes$ is the Kronecker (Zehfuss) product for matrices and the outer product for vectors.

\end{frame}
